Os resultados revelam padrões importantes sobre code review em repositórios
populares do GitHub.

\textbf{Dimensão A -- Status do PR:}
O tempo de análise é o único fator com correlação moderada ao status de merge
(RQ02, $\rho = -0,3910$). Isso sugere que PRs que são integrados rapidamente
têm maior chance de merge. Possíveis explicações incluem: PRs simples e bem
preparados passam mais rápido pela revisão, ou PRs que demoram muito perdem
relevância ou conflitam com mudanças mais recentes.

As demais métricas (tamanho, descrição, participantes) apresentam correlações
desprezíveis com o status.

\textbf{Dimensão B -- Número de Revisões:}
O tamanho do PR (RQ05, $\rho = 0,3114$) mostra correlação moderada com o
número de revisões. PRs maiores demandam mais atenção de revisores, como
evidenciado pelo crescimento da mediana de alterações de 14 (1 revisão)
para 342 (10 revisões). O tempo de análise (RQ06, $\rho = 0,0917$) e
o tamanho da descrição (RQ07, $\rho = 0,1861$) apresentam correlações
desprezíveis e fracas, respectivamente.

\textbf{Limitações:}
Os campos de comentários não foram corretamente coletados, limitando nossa
análise de ``interações'' ao número de participantes. Além disso, a amostra
compreende apenas repositórios populares, o que pode não refletir projetos
menores ou corporativos. A winsorização no percentil 99, embora necessária
para mitigar outliers extremos, pode ter removido informações relevantes
sobre PRs excepcionalmente grandes.
